\chapter{Metodologia}

\section{Abordagem de Pesquisa}

Este trabalho adota a metodologia de \textit{Experience Report} (Relato de Experiência) para documentar o processo de concepção, desenvolvimento e implementação da arquitetura Data Fabric DermAlert. Essa abordagem é especialmente adequada para pesquisas que visam compartilhar aprendizados práticos advindos da implementação de soluções tecnológicas inovadoras em contextos reais, permitindo a análise crítica de desafios, decisões arquiteturais e lições aprendidas ao longo do desenvolvimento. A escolha por este método se fundamenta na natureza aplicada do projeto, que integra aspectos tecnológicos, regulatórios e clínicos na criação de uma solução para integração de dados dermatológicos heterogêneos no âmbito do Sistema Único de Saúde (SUS) brasileiro.

\section{Perguntas de Pesquisa}

O desenvolvimento da arquitetura DermAlert foi orientado pelas seguintes perguntas de pesquisa:

\begin{enumerate}
    \item Como implementar uma arquitetura Data Fabric que atenda aos requisitos específicos de armazenamento, processamento e governança de imagens dermatológicas e metadados clínicos associados?
    \item Quais são os principais desafios técnicos, operacionais e regulamentares na implementação de uma arquitetura Data Fabric para dados médicos sensíveis em conformidade com a LGPD?
    \item De que forma a integração semântica pode resolver a heterogeneidade entre diferentes instituições de saúde em um sistema distribuído?
    \item Como a virtualização e federação de dados podem melhorar a eficiência no acesso a informações distribuídas sem comprometer a segurança e governança?
    \item Qual o impacto da arquitetura proposta na qualidade, diversidade e acessibilidade dos dados para o desenvolvimento de modelos de \textit{machine learning} mais equitativos?
\end{enumerate}

\section{Processo de Desenvolvimento da Arquitetura}

O desenvolvimento da arquitetura foi dividido nas seguintes fases:

\subsection{Fase 1: Levantamento de Requisitos e Análise de Contexto}

O processo iniciou-se com uma análise abrangente do contexto dermatológico brasileiro, incluindo:
\begin{itemize}
    \item Desafios epidemiológicos: análise de estatísticas do INCA sobre câncer de pele e necessidades de detecção precoce;
    \item Limitações tecnológicas: identificação de vieses em bases de dados existentes (como o ISIC Archive) e necessidade de maior diversidade de tons de pele;
    \item Requisitos regulatórios: mapeamento das exigências da LGPD e normas de segurança em saúde;
    \item Necessidades operacionais: levantamento de fluxos de trabalho e limitações de infraestrutura junto a profissionais da Atenção Primária.
\end{itemize}

\subsection{Fase 2: Concepção Arquitetural}

A arquitetura foi concebida com base nos seguintes princípios:
\begin{itemize}
    \item Descentralização: manutenção dos dados em suas fontes originais;
    \item Elasticidade e escalabilidade: expansão dinâmica da infraestrutura;
    \item Segurança e governança: conformidade regulatória desde o design;
    \item Interoperabilidade: integração semântica entre fontes heterogêneas;
    \item Custo-efetividade: uso de tecnologias \textit{open source}.
\end{itemize}

\subsection{Fase 3: Seleção e Integração Tecnológica}

A seleção das tecnologias considerou cinco critérios principais:
\begin{itemize}
    \item Escalabilidade;
    \item Manutenibilidade;
    \item Resiliência;
    \item Custo-efetividade;
    \item Flexibilidade.
\end{itemize}

O \textit{stack} tecnológico implementado incluiu:
\begin{itemize}
    \item Armazenamento distribuído: MinIO (arquitetura em camadas: bronze, prata, ouro);
    \item Gerenciamento transacional: Delta Lake para metadados, versionamento e auditoria;
    \item Orquestração: Apache Airflow para automação de pipelines;
    \item Processamento: Apache Spark para análises em larga escala;
    \item Compartilhamento: Delta Sharing para colaboração segura;
    \item API: FastAPI para interface de acesso;
    \item Containerização: Kubernetes para alta disponibilidade;
    \item Governança: catálogo integrado de metadados com IA/ML para descoberta automática.
    \item Interface de usuário: Next.js/React com Redux para desenvolvimento do frontend, proporcionando experiências reativas, escaláveis e integradas ao backend.
\end{itemize}

\subsection{Fase 4: Implementação da Camada de Integração Semântica}

Desenvolvimento de sistema de equivalências terminológicas, incluindo:
\begin{itemize}
    \item Mapeamento direto e agrupamentos hierárquicos de termos clínicos;
    \item Propagação transitiva de equivalências na rede semântica;
    \item Interoperabilidade multi-institucional entre diferentes nomenclaturas.
\end{itemize}

\subsection{Fase 5: Desenvolvimento do Aplicativo de Coleta}

Desenvolvimento do aplicativo DermAlert para uso em postos de saúde, contemplando:
\begin{itemize}
    \item Fluxo clínico otimizado;
    \item Captura padronizada de imagens;
    \item Consentimento digital integrado;
    \item Anonimização dos dados coletados;
    \item Sincronização assíncrona com o Data Fabric.
\end{itemize}

\section{Estratégia de Validação}

\subsection{Validação Técnica}

\begin{itemize}
    \item Testes de integração entre componentes e consistência do pipeline;
    \item Avaliação de desempenho: latência de consultas, throughput e escalabilidade;
    \item Segurança e conformidade: auditoria de criptografia, anonimização, RBAC e aderência à LGPD.
\end{itemize}

\subsection{Validação Operacional}

\begin{itemize}
    \item Testes com usuários reais: avaliação de usabilidade e aderência aos fluxos clínicos;
    \item Piloto controlado em cinco unidades de saúde: monitoramento de métricas operacionais e impacto no trabalho.
\end{itemize}

\section{Critérios de Avaliação}

\subsection{Métricas Quantitativas}

\begin{itemize}
    \item Escalabilidade: suporte a uso simultâneo sem degradação;
    \item Disponibilidade: uptime superior a 99,5\%;
    \item Latência: resposta inferior a 2s para consultas padrão;
    \item Throughput: processamento de até 1000 imagens/hora;
    \item Qualidade de dados: métricas para imagens (resolução, SNR, nitidez, contraste).
\end{itemize}

\subsection{Métricas Qualitativas}

\begin{itemize}
    \item Usabilidade: questionários estruturados com profissionais;
    \item Interoperabilidade: integração com fontes heterogêneas;
    \item Governança: eficácia de rastreabilidade e auditoria;
    \item Diversidade: representatividade de tons de pele na base.
\end{itemize}



\section{Considerações sobre Sustentabilidade}

\subsection{Sustentabilidade Tecnológica}

\begin{itemize}
    \item Open source: independência de fornecedores;
    \item Modularidade: evolução e substituição de componentes;
    \item Padronização: aderência a padrões internacionais de interoperabilidade;
    \item Comunidade: engajamento com desenvolvedores.
\end{itemize}

\subsection{Sustentabilidade Operacional}

\begin{itemize}
    \item Capacitação de equipes técnicas;
    \item Documentação operacional e de manutenção;
    \item Suporte técnico continuado;
    \item Roadmap de evolução futura.
\end{itemize}

\section{Considerações Finais sobre a Metodologia}

A abordagem adotada proporcionou base sólida para o desenvolvimento da arquitetura DermAlert, garantindo não só viabilidade técnica, mas também adequação ao contexto da saúde pública brasileira, com foco em equidade, segurança e sustentabilidade.